
\begin{register}{H}{mvendorid}{0xF11}\label{regdesc:MVENDORID}
\regfield{MVENDORID}{32}{0}{{0x00000612}}%
\reglabel{Reset}\regnewline%
\begin{regdesc}\begin{reglist}
\label{fielddesc:MVENDORID}\item [MVENDORID] 供应商编号。（只读）
\end{reglist}\end{regdesc}
\end{register}

\begin{register}{H}{marchid}{0xF12}\label{regdesc:MARCHID}
\regfield{MARCHID}{32}{0}{{0x80000001}}%
\reglabel{Reset}\regnewline%
\begin{regdesc}\begin{reglist}
\label{fielddesc:MARCHID}\item [MARCHID] 架构编号。（只读）
\end{reglist}\end{regdesc}
\end{register}

\begin{register}{H}{mimpid}{0xF13}\label{regdesc:MIMPID}
\regfield{MIMPID}{32}{0}{{0x00000002}}%
\reglabel{Reset}\regnewline%
\begin{regdesc}\begin{reglist}
\label{fielddesc:MIMPID}\item [MIMPID] 实现编号。（只读）
\end{reglist}\end{regdesc}
\end{register}

\begin{register}{H}{mhartid}{0xF14}\label{regdesc:MHARTID}
\regfield{MHARTID}{32}{0}{{0x00000000}}%
\reglabel{Reset}\regnewline%
\begin{regdesc}\begin{reglist}
\label{fielddesc:MHARTID}\item [MHARTID] 硬件线程编号。（只读）
\end{reglist}\end{regdesc}
\end{register}

\begin{register}{H}{mstatus}{0x300}\label{regdesc:MSTATUS}
\regfield{(reserved)}{10}{22}{{0x000}}%
\regfield{TW}{1}{21}{{0}}%
\regfield{(reserved)}{8}{13}{{0x00}}%
\regfield{MPP}{2}{11}{{0x0}}%
\regfield{(reserved)}{3}{8}{{0x0}}%
\regfield{MPIE}{1}{7}{{0}}%
\regfield{(reserved)}{3}{4}{{0x0}}%
\regfield{MIE}{1}{3}{{0}}%
\regfield{(reserved)}{3}{0}{{0x0}}%
\reglabel{Reset}\regnewline%
\begin{regdesc}\begin{reglist}
\label{fielddesc:MSTATUSMIE}\item [MIE] 全局机器模式中断使能。（读/写）
\label{fielddesc:MSTATUSMPIE}\item [MPIE] 之前的 \hyperref[fielddesc:MSTATUSMIE]{MIE}。（读/写）
\label{fielddesc:MSTATUSMPP}\item [MPP] 机器之前的特权模式。（读/写）
\newline 可能的值：
\begin{itemize}
    \setlength\itemsep{-1em}
    \item 0x0：用户模式
    \item 0x3：机器模式
\end{itemize}
\textit{说明：仅低位可写。由于高位直接绑定低位，写入高位将被忽略。}
\label{fielddesc:MSTATUSTW}\item [TW] 超时等待。（读/写）
\newline 如果该位置 1，用户模式下的 WFI（等待中断）指令将导致非法指令异常。
\end{reglist}\end{regdesc}
\end{register}

\begin{register}{H}{misa}{0x301}\label{regdesc:MISA}
\regfield{MXL}{2}{30}{{0x1}}%
\regfield{(reserved)}{4}{26}{{0x0}}%
\regfield{Z}{1}{25}{{0}}%
\regfield{Y}{1}{24}{{0}}%
\regfield{X}{1}{23}{{0}}%
\regfield{W}{1}{22}{{0}}%
\regfield{V}{1}{21}{{0}}%
\regfield{U}{1}{20}{{1}}%
\regfield{T}{1}{19}{{0}}%
\regfield{S}{1}{18}{{0}}%
\regfield{R}{1}{17}{{0}}%
\regfield{Q}{1}{16}{{0}}%
\regfield{P}{1}{15}{{0}}%
\regfield{O}{1}{14}{{0}}%
\regfield{N}{1}{13}{{0}}%
\regfield{M}{1}{12}{{1}}%
\regfield{L}{1}{11}{{0}}%
\regfield{K}{1}{10}{{0}}%
\regfield{J}{1}{9}{{0}}%
\regfield{I}{1}{8}{{1}}%
\regfield{H}{1}{7}{{0}}%
\regfield{G}{1}{6}{{0}}%
\regfield{F}{1}{5}{{0}}%
\regfield{E}{1}{4}{{0}}%
\regfield{D}{1}{3}{{0}}%
\regfield{C}{1}{2}{{1}}%
\regfield{B}{1}{1}{{0}}%
\regfield{A}{1}{0}{{0}}%
\reglabel{Reset}\regnewline%
\begin{regdesc}\begin{reglist}
\label{fielddesc:MISAMXL}\item [MXL] 机器 XLEN = 1（32 位）。（只读）
\label{fielddesc:MISAZ}\item [Z] 保留 = 0。（只读）
\label{fielddesc:MISAY}\item [Y] 保留 = 0。（只读）
\label{fielddesc:MISAX}\item [X] 非标准扩展 = 0。（只读）
\label{fielddesc:MISAW}\item [W] 保留 = 0。（只读）
\label{fielddesc:MISAV}\item [V] 保留 = 0。（只读）
\label{fielddesc:MISAU}\item [U] 实现用户模式 = 1。（只读）
\label{fielddesc:MISAT}\item [T] 保留 = 0。（只读）
\label{fielddesc:MISAS}\item [S] 实现监督模式 = 0。（只读）
\label{fielddesc:MISAR}\item [R] 保留 = 0。（只读）
\label{fielddesc:MISAQ}\item [Q] 四精度浮点扩展 = 0。（只读）
\label{fielddesc:MISAP}\item [P] 保留 = 0。（只读）
\label{fielddesc:MISAO}\item [O] 保留 = 0。（只读）
\label{fielddesc:MISAN}\item [N] 支持用户级别中断 = 0。（只读）
\label{fielddesc:MISAM}\item [M] 整数乘除法标准扩展 = 1。（只读）
\label{fielddesc:MISAL}\item [L] 保留 = 0。（只读）
\label{fielddesc:MISAK}\item [K] 保留 = 0。（只读）
\label{fielddesc:MISAJ}\item [J] 保留 = 0。（只读）
\label{fielddesc:MISAI}\item [I] RV32I 基本 ISA = 1。（只读）
\label{fielddesc:MISAH}\item [H] 虚拟机管理程序扩展 = 0。（只读）
\label{fielddesc:MISAG}\item [G] 其他标准扩展 = 0。（只读）
\label{fielddesc:MISAF}\item [F] 单精度浮点扩展 = 0。（只读）
\label{fielddesc:MISAE}\item [E] RV32E 基本 ISA = 0。（只读）
\label{fielddesc:MISAD}\item [D] 双精度浮点扩展 = 0。（只读）
\label{fielddesc:MISAC}\item [C] 压缩标准扩展 = 1。（只读）
\label{fielddesc:MISAB}\item [B] 保留 = 0。（只读）
\label{fielddesc:MISAA}\item [A] 原子标准扩展 = 0。（只读）
\end{reglist}\end{regdesc}
\end{register}

\begin{register}{H}{mtvec}{0x305}\label{regdesc:MTVEC}
\regfield{BASE}{24}{8}{{0x000000}}%
\regfield{(reserved)}{6}{2}{{0x00}}%
\regfield{MODE}{2}{0}{{0x1}}%
\reglabel{Reset}\regnewline%
\begin{regdesc}\begin{reglist}
\label{fielddesc:MTVECMODE}\item [MODE] 仅支持向量模式 \textbf{0x1}。（只读）
\label{fielddesc:MTVECBASE}\item [BASE] 异常向量基址的高 24 位为 256 字节对齐。（读/写）
\end{reglist}\end{regdesc}
\end{register}

\begin{register}{H}{mscratch}{0x340}\label{regdesc:MSCRATCH}
\regfield{MSCRATCH}{32}{0}{{0x00000000}}%
\reglabel{Reset}\regnewline%
\begin{regdesc}\begin{reglist}
\label{fielddesc:MSCRATCH}\item [MSCRATCH] 用户自定义的机器暂存寄存器。（读/写）
\end{reglist}\end{regdesc}
\end{register}

\begin{register}{H}{mepc}{0x341}\label{regdesc:MEPC}
\regfield{MEPC}{32}{0}{{0x00000000}}%
\reglabel{Reset}\regnewline%
\begin{regdesc}\begin{reglist}
\label{fielddesc:MEPC}\item [MEPC] 机器陷阱/异常程序计数器。（读/写） \newline 当 CPU 遇到异常时，此域将自动更新为 CPU 将要执行的指令的地址。
\end{reglist}\end{regdesc}
\end{register}

\begin{register}{H}{mcause}{0x342}\label{regdesc:MCAUSE}
\regfield{Interrupt Flag}{1}{31}{{0}}%
\regfield{(reserved)}{26}{5}{{0x0000000}}%
\regfield{Exception Code}{5}{0}{{0x00}}%
\reglabel{Reset}\regnewline%
\begin{regdesc}\begin{reglist}
\label{fielddesc:MCAUSECODE}\item [Exception Code] CPU 进入异常时，此域将自动更新为最近的异常或中断的唯一 ID。（读/写）
\newline 可能的异常 ID：
\begin{itemize}
    \setlength\itemsep{-1em}
    \item 0x1：PMP 指令访问错误
    \item 0x2：非法指令
    \item 0x3：硬件断点/观察点或 EBREAK
    \item 0x5：PMP 读存储器访问错误
    \item 0x7：PMP 写存储器访问错误
    \item 0x8：用户模式环境调用 (ECALL)
    \item 0xb：机器模式环境调用
\end{itemize}
\textit {说明：异常 ID 0x0（指令地址非对齐）不存在，因为 CPU 在取指时始终屏蔽地址的最低位。}
\label{fielddesc:MCAUSEINT}\item [Interrupt Flag] CPU 进入异常时，此标志位将自动更新。（读/写） \newline 如果被置位，则表示最近的陷阱是由中断引起。在异常情况下保持为 0。 \newline \textit{说明：中断控制器将中断编号 1-31 全部用于外部中断源，而 RISC-V 标准则为内核的内部中断源预留了编号 0-15。}
\end{reglist}\end{regdesc}
\end{register}

\begin{register}{H}{mtval}{0x343}\label{regdesc:MTVAL}
\regfield{MTVAL}{32}{0}{{0x00000000}}%
\reglabel{Reset}\regnewline%
\begin{regdesc}\begin{reglist}
\label{fielddesc:MTVAL}\item [MTVAL] 机器模式异常值。（读/写） \newline 将自动更新为与异常有关的数据，该数据可能有助于处理该异常。
\newline 根据异常编号有以下解读：
\begin{itemize}
    \setlength\itemsep{-1em}
    \item 0x1：指令虚拟地址错误
    \item 0x2：指令 opcode 错误
    \item 0x5：存储器读操作的数据地址错误
    \item 0x7：存储器写操作的数据地址错误
\end{itemize}
\textit {说明：该寄存器不支持其他异常 ID 和中断。}
\end{reglist}\end{regdesc}
\end{register}

\begin{register}{H}{mpcer}{0x7E0}\label{regdesc:MPCER}
\regfield{(reserved)}{21}{11}{{0x000}}%
\regfield{INST\_COMP}{1}{10}{{0}}%
\regfield{(BRANCH\_TAKEN}{1}{9}{{0}}%
\regfield{BRANCH}{1}{8}{{0}}%
\regfield{JMP\_UNCOND}{1}{7}{{0}}%
\regfield{STORE}{1}{6}{{0}}%
\regfield{LOAD}{1}{5}{{0}}%
\regfield{IDLE}{1}{4}{{0}}%
\regfield{JMP\_HAZARD}{1}{3}{{0}}%
\regfield{LD\_HAZARD}{1}{2}{{0}}%
\regfield{INST}{1}{1}{{0}}%
\regfield{CYCLE}{1}{0}{{0}}%
\reglabel{Reset}\regnewline%
\begin{regdesc}\begin{reglist}
\label{fielddesc:INST_COMP}\item [INST\_COMP] 计数压缩指令。（读/写）
\label{fielddesc:BRANCH_TAKEN}\item [BRANCH\_TAKEN] 计数跳转的分支。（读/写）
\label{fielddesc:BRANCH}\item [BRANCH] 计数分支。（读/写）
\label{fielddesc:JUMP_UNCONDITIONAL}\item [JMP\_UNCOND] 计数无条件跳转。（读/写）
\label{fielddesc:STORE}\item [STORE] 计数存储器写操作。（读/写）
\label{fielddesc:LOAD}\item  [LOAD] 计数存储器读操作。（读/写）
\label{fielddesc:IDLE}\item  [IDLE] 计数 IDLE 周期。（读/写）
\label{fielddesc:JUMP_HAZARDS}\item  [JMP\_HAZARD] 计数跳转冲突。（读/写）
\label{fielddesc:LOAD_HAZARDS}\item  [LD\_HAZARD] 计数存储器读操作冲突。（读/写）
\label{fielddesc:INSTRUCTION}\item  [INST] 计数指令。（读/写）
\label{fielddesc:CYCLES}\item  [CYCLE] 计数时钟周期，WFI 模式下周期计数不增加。（读/写）

\textit{注意：每个位选择一个特定事件由计数器递增计数。如果多个事件被选择且同时发生，则计数器只递增 1。 }

\end{reglist}\end{regdesc}
\end{register}

\begin{register}{H}{mpcmr}{0x7E1}\label{regdesc:MPCMR}
\regfield{(reserved)}{30}{2}{{0}}%
\regfield{COUNT\_SAT}{1}{1}{{1}}%
\regfield{COUNT\_EN}{1}{0}{{1}}%
\reglabel{Reset}\regnewline%
\begin{regdesc}\begin{reglist}
\label{fielddesc:COUNTER_SATURATION}\item [COUNT\_SAT] 计数器饱和控制。（读/写）
\newline 可能的值：
\begin{itemize}
    \setlength\itemsep{-1em}
    \item 0：超出最大值上溢
    \item 1：超出最大值暂停
\end{itemize}
\label{fielddesc:COUNTER_ENABLE}\item [COUNT\_EN] 计数器使能控制。（读/写）
\newline 可能的值：
\begin{itemize}
    \setlength\itemsep{-1em}
    \item 0：禁能
    \item 1：使能
\end{itemize}
\end{reglist}\end{regdesc}
\end{register}

\begin{register}{H}{mpccr}{0x7E2}\label{regdesc:MPCCR}
\regfield{MPCCR}{32}{0}{{0x00000000}}%
\reglabel{Reset}\regnewline%
\begin{regdesc}\begin{reglist}
\label{fielddesc:MPCCR}\item [MPCCR] 性能计数器计数的值。（读/写）
\end{reglist}\end{regdesc}
\end{register}

\begin{register}{H}{cpu\_gpio\_oen}{0x803}\label{regdesc:GPIO_OEN}
\regfield{(reserved)}{24}{8}{{0}}%
\regfield{CPU\_GPIO\_OEN[7]}{1}{7}{{0}}%
\regfield{CPU\_GPIO\_OEN[6]}{1}{6}{{0}}%
\regfield{CPU\_GPIO\_OEN[5]}{1}{5}{{0}}%
\regfield{CPU\_GPIO\_OEN[4]}{1}{4}{{0}}%
\regfield{CPU\_GPIO\_OEN[3]}{1}{3}{{0}}%
\regfield{CPU\_GPIO\_OEN[2]}{1}{2}{{0}}%
\regfield{CPU\_GPIO\_OEN[1]}{1}{1}{{0}}%
\regfield{CPU\_GPIO\_OEN[0]}{1}{0}{{0}}%
\reglabel{Reset}\regnewline%
\begin{regdesc}\begin{reglist}
\label{fielddesc:GPIO_OEN}\item [CPU\_GPIO\_OEN] GPIOn (n=0 \verb+~+ 21) 输出使能。CPU\_GPIO\_OEN[7:0] 分别对应章节 \textit{\nameref{mod:iomuxgpio}} 中表 \ref{tab:iomuxgpio-periph-signals-via-gpio-matrix} 里的 cpu\_gpio\_out\_oen[7:0] 输出使能信号。CPU\_GPIO\_OEN 的值与 cpu\_gpio\_out\_oen 的值对应。
\newline 此寄存器是 \hyperref[fielddesc:GPIO_OUT]{CPU\_GPIO\_OUT} 的使能寄存器。（读/写）

\begin{itemize}
    \setlength\itemsep{-1em}
    \item 0：GPIO 输出关闭
    \item 1：GPIO 输出使能
\end{itemize}
\end{reglist}\end{regdesc}
\end{register}

\begin{register}{H}{cpu\_gpio\_in}{0x804}\label{regdesc:GPIO_IN}
\regfield{(reserved)}{24}{8}{{0}}%
\regfield{CPU\_GPIO\_IN[7]}{1}{7}{{0}}%
\regfield{CPU\_GPIO\_IN[6]}{1}{6}{{0}}%
\regfield{CPU\_GPIO\_IN[5]}{1}{5}{{0}}%
\regfield{CPU\_GPIO\_IN[4]}{1}{4}{{0}}%
\regfield{CPU\_GPIO\_IN[3]}{1}{3}{{0}}%
\regfield{CPU\_GPIO\_IN[2]}{1}{2}{{0}}%
\regfield{CPU\_GPIO\_IN[1]}{1}{1}{{0}}%
\regfield{CPU\_GPIO\_IN[0]}{1}{0}{{0}}%
\reglabel{Reset}\regnewline%
\begin{regdesc}\begin{reglist}
\label{fielddesc:GPIO_IN}\item [CPU\_GPIO\_IN] 读取 SoC GPIOn (n=0 \verb+~+ 21) 的输入值（1 为高电平，0 为低电平）。
\newline CPU\_GPIO\_IN[7:0] 分别对应章节 \textit{\nameref{mod:iomuxgpio}} 中表 \ref{tab:iomuxgpio-periph-signals-via-gpio-matrix} 里的 cpu\_gpio\_in[7:0] 输入信号。
\newline CPU\_GPIO\_IN[7:0] 只能通过 GPIO 交换矩阵映射到 GPIO。详细描述请参考章节 \ref{sec:peripheral-input-via-gpio-matrix}。（只读）
\end{reglist}\end{regdesc}
\end{register}


\begin{register}{H}{cpu\_gpio\_out}{0x805}\label{regdesc:GPIO_OUT}
\regfield{(reserved)}{24}{8}{{0}}%
\regfield{CPU\_GPIO\_OUT[7]}{1}{7}{{0}}%
\regfield{CPU\_GPIO\_OUT[6]}{1}{6}{{0}}%
\regfield{CPU\_GPIO\_OUT[5]}{1}{5}{{0}}%
\regfield{CPU\_GPIO\_OUT[4]}{1}{4}{{0}}%
\regfield{CPU\_GPIO\_OUT[3]}{1}{3}{{0}}%
\regfield{CPU\_GPIO\_OUT[2]}{1}{2}{{0}}%
\regfield{CPU\_GPIO\_OUT[1]}{1}{1}{{0}}%
\regfield{CPU\_GPIO\_OUT[0]}{1}{0}{{0}}%
\reglabel{Reset}\regnewline%
\begin{regdesc}\begin{reglist}
\label{fielddesc:GPIO_OUT}\item [CPU\_GPIO\_OUT] 向 SoC GPIOn (n=0 \verb+~+ 21) 写输出值（1 为高电平，0 为低电平）。\hyperref[fielddesc:GPIO_OEN]{CPU\_GPIO\_OEN} 置位时，写输出值才有效。
\newline CPU\_GPIO\_OUT[7:0] 分别对应章节 \textit{\nameref{mod:iomuxgpio}} 中表 \ref{tab:iomuxgpio-periph-signals-via-gpio-matrix} 里的 cpu\_gpio\_out[7:0] 输出信号。
\newline CPU\_GPIO\_OUT[7:0] 只能通过 GPIO 交换矩阵映射到 GPIO。详细描述请参考章节 \ref{sec:peripheral-output-via-gpio-matrix}。（读/写）
\end{reglist}\end{regdesc}
\end{register}
